\documentclass[11pt, a4paper]{article}
\usepackage[margin=2.5cm]{geometry}
\usepackage{amsmath, amssymb, amsthm}
\usepackage{booktabs}
\usepackage{hyperref}
\usepackage{microtype}
\usepackage{parskip}

\newcommand{\E}{\mathbb{E}}
\newcommand{\Var}{\mathrm{Var}}
\newcommand{\Cov}{\mathrm{Cov}}
\newcommand{\R}{\mathbb{R}}
\newcommand{\xt}{\tilde{x}}
\newcommand{\Xt}{\widetilde{X}}
\newcommand{\hbeta}{\hat{\beta}}
\newcommand{\halpha}{\hat{\alpha}}
\newcommand{\htau}{\hat{\tau}}
\newcommand{\hsigma}{\hat{\sigma}}

\title{Variance Estimation in \texttt{dsFedAteClient}\\[0.4em]
       \large Continuous vs.\ Binary Outcomes, Pooled vs.\ Arm-Specific~$\sigma^2$}
\author{dsFedAteClient package notes}
\date{}

\begin{document}
\maketitle
\tableofcontents
\bigskip

% -----------------------------------------------------------------------
\section{Setup and notation}

We observe $n = \sum_{k=1}^K n_k$ units split across $K$ sites.
At each site $k$, unit $i$ has covariates $x_i \in \R^d$, a binary treatment
indicator $W_i \in \{0,1\}$, and an outcome $Y_i$.

We fit two arm-specific models (one per treatment value) using federated GLM,
broadcasting the pooled coefficients
$(\halpha_1,\hbeta_1)$ (treated) and $(\halpha_0,\hbeta_0)$ (control)
to all sites.  Write $\xt_i = (1, x_i^\top)^\top \in \R^{d+1}$ for the
augmented covariate vector.

% -----------------------------------------------------------------------
\section{The federated ATE estimator}

The global ATE estimator is the sample-size-weighted average of local estimates:
\begin{equation}\label{eq:ate}
  \htau = \frac{1}{n}\sum_{k=1}^K n_k \htau_k.
\end{equation}

\paragraph{Continuous outcome (OLS).}
Each arm is fitted with a linear model $Y_i = \alpha_w + x_i^\top\beta_w + \varepsilon_i$.
The local ATE is
\begin{equation}\label{eq:tau_cont}
  \htau_k = (\halpha_1 - \halpha_0) + \bar{x}_k^\top(\hbeta_1 - \hbeta_0),
\end{equation}
where $\bar{x}_k = n_k^{-1}\sum_{i \in k} x_i$ is the local covariate mean.

\paragraph{Binary outcome (logistic, G-computation).}
Each arm is fitted with a logistic model
$\mu_w(x_i) = \sigma(\halpha_w + x_i^\top\hbeta_w)$,
where $\sigma(\cdot)$ is the logistic function.
The local ATE is
\begin{equation}\label{eq:tau_bin}
  \htau_k = \frac{1}{n_k}\sum_{i \in k}\bigl[\mu_1(x_i) - \mu_0(x_i)\bigr].
\end{equation}

% -----------------------------------------------------------------------
\section{Delta-method variance: unified formula}

By the delta method applied to the mapping
$(\halpha_1,\hbeta_1,\halpha_0,\hbeta_0) \mapsto \htau$,
the asymptotic variance of $\htau$ decomposes into a
\emph{parametric} term (uncertainty in the GLM coefficients) and a
\emph{sampling} term (finite-sample variability of the local ATEs):
\begin{equation}\label{eq:var}
  \Var(\htau)
    = \underbrace{g_1^\top V_1\, g_1 + g_0^\top V_0\, g_0}_{\text{parametric}}
    + \underbrace{\frac{1}{n(n-1)}\sum_{k=1}^K\Bigl[
        s_k + n_k\bigl(\htau_k - \htau\bigr)^2\Bigr]}_{\text{sampling}},
\end{equation}
where:
\begin{itemize}
  \item $V_w \in \R^{(d+1)\times(d+1)}$ is the covariance matrix of
        $(\halpha_w,\hbeta_w^\top)^\top$,
  \item $g_1, g_0 \in \R^{d+1}$ are gradient vectors (defined per outcome type
        in Sections~\ref{sec:cont} and~\ref{sec:bin}),
  \item $s_k = \sum_{i \in k}(\delta_i - \htau_k)^2$ with
        $\delta_i = \hat\mu_1(x_i) - \hat\mu_0(x_i)$ is the local scatter of
        individual treatment effect estimates around the local mean.
\end{itemize}

This single formula covers both outcome types.  The differences between
continuous and binary lie entirely in how $g_w$ and $V_w$ are defined.

% -----------------------------------------------------------------------
\section{Continuous outcomes}\label{sec:cont}

\subsection{Gradient vectors}

Under the linear model, $\delta_i = \xt_i^\top(\halpha_1-\halpha_0,\hbeta_1^\top-\hbeta_0^\top)^\top$
is linear in the coefficients, so the gradient of $\htau$ with respect to
$(\halpha_1,\hbeta_1^\top)^\top$ evaluated at the pooled fit is
\begin{equation}
  g_1 = \frac{1}{n}\sum_{k=1}^K\sum_{i \in k}\xt_i
       = \bigl(1,\, \bar{x}^\top\bigr)^\top =: \bar{\xt},
\end{equation}
and $g_0 = -\bar{\xt}$ (same magnitude, opposite sign because increasing the
control-arm intercept decreases the ATE).

In practice $\bar{\xt}$ is assembled from per-site raw sums returned by
\texttt{localAteDS}: $g_{1,\text{raw},k} = g_{0,\text{raw},k} = \sum_{i \in k}\xt_i$,
and the client forms $g_1 = \sum_k g_{1,\text{raw},k}/n$.

\subsection{The covariance matrix $V_w$}

For OLS the coefficient covariance is
\begin{equation}\label{eq:V_ols}
  V_w = \sigma^2 \bigl(X_w^\top X_w\bigr)^{-1},
\end{equation}
where $X_w$ is the $n_w \times (d+1)$ design matrix of the $w$-arm and
$\sigma^2 = \Var(\varepsilon_i)$ is the noise variance, assumed the same across
both arms and all sites.

\subsection{Estimating $\sigma^2$ via pooled residual sums}\label{sec:pooled}

\texttt{ds.glm} with \texttt{viewVarCov = TRUE} returns the
\emph{unscaled} matrix $(X_w^\top X_w)^{-1}$, not the full covariance~\eqref{eq:V_ols}.
We therefore estimate $\sigma^2$ separately from the residual sums of squares
collected from all sites and both arms:
\begin{equation}\label{eq:sigma2}
  \hsigma^2
    = \frac{\displaystyle\sum_{k=1}^K \mathrm{RSS}_k^{(1)} + \mathrm{RSS}_k^{(0)}}
           {n - 2(d+1)},
\end{equation}
where $\mathrm{RSS}_k^{(w)} = \sum_{i \in k,\, W_i=w}(Y_i - \halpha_w - x_i^\top\hbeta_w)^2$.
The denominator accounts for $2(d+1)$ estimated parameters (intercept plus $d$
slopes for each of the two arms).

The full covariance used in~\eqref{eq:var} is then
$V_w = \hsigma^2 (X_w^\top X_w)^{-1}$,
with $(X_w^\top X_w)^{-1}$ taken from \texttt{ds.glm\$VarCovMatrix}.

\subsection{Comparison with arm-specific $\sigma^2$}

An alternative is to use the arm-specific estimate
$\hsigma_w^2 = \mathrm{RSS}^{(w)}/(n_w - (d+1))$,
which is what standard \texttt{glm} embeds in \texttt{vcov(fit)}.
The pooled estimator~\eqref{eq:sigma2} is preferable when:
\begin{enumerate}
  \item \textbf{Homoscedasticity across arms} holds (single $\sigma^2$ for both
        arms).  In an RCT this is the natural assumption when treatment does not
        affect the error distribution.
  \item \textbf{Efficiency}: the pooled estimate has $n - 2(d+1)$ degrees of
        freedom versus $n_w - (d+1)$ for each arm separately, halving the
        estimation noise when arms are balanced.
\end{enumerate}
If treatment shifts the variance of $Y$, arm-specific estimates are more robust,
but this situation lies outside the standard RCT model assumed throughout.

% -----------------------------------------------------------------------
\section{Binary outcomes}\label{sec:bin}

\subsection{Gradient vectors}

Under the logistic model, $\delta_i = \mu_1(x_i) - \mu_0(x_i)$ is non-linear
in the coefficients.  The gradient of $\htau$ with respect to
$(\halpha_1,\hbeta_1^\top)^\top$ is
\begin{equation}
  g_1 = \frac{1}{n}\sum_{k=1}^K\sum_{i \in k} v_1(x_i)\,\xt_i,
  \qquad
  v_w(x_i) = \mu_w(x_i)\bigl[1 - \mu_w(x_i)\bigr],
\end{equation}
where $v_w(x_i)$ is the logistic variance weight (the derivative of $\sigma$).
Analogously, $g_0 = -n^{-1}\sum_{k,i} v_0(x_i)\,\xt_i$.

In practice \texttt{localAteDS} returns per-site raw sums
$g_{1,\text{raw},k} = \sum_{i \in k} v_1(x_i)\xt_i$ and
$g_{0,\text{raw},k} = \sum_{i \in k} v_0(x_i)\xt_i$;
the client assembles $g_1 = \sum_k g_{1,\text{raw},k}/n$ (and similarly $g_0$
with a sign flip).

\subsection{The covariance matrix $V_w$ and why there is no $\sigma^2$}

For logistic regression the observed Fisher information is
\begin{equation}
  H_w = X_w^\top\,\mathrm{diag}\bigl(v_w(x_i)\bigr)\,X_w,
\end{equation}
and the maximum-likelihood covariance of $(\halpha_w,\hbeta_w^\top)^\top$ is
\begin{equation}\label{eq:V_bin}
  V_w = H_w^{-1}.
\end{equation}
The key difference from OLS is that $H_w$ already encodes the observation-level
variance $\Var(Y_i \mid x_i) = \mu_w(x_i)(1-\mu_w(x_i))$; there is no
separate scale parameter $\sigma^2$ to estimate.  Formally, the binomial family
has \emph{dispersion fixed at~1}.

This is why \texttt{ds.glm(viewVarCov = TRUE)\$VarCovMatrix} for a binomial
model returns the \emph{scaled} covariance $H_w^{-1}$ directly (dispersion = 1,
so no scaling is needed), whereas for Gaussian it returns the unscaled
$(X_w^\top X_w)^{-1}$.
In the implementation this is handled uniformly by setting
$\hsigma^2 = 1$ for binary and multiplying:
\begin{equation}
  V_w = \hsigma^2 \times \texttt{VarCovMatrix}
      = \begin{cases}
          \hsigma^2 (X_w^\top X_w)^{-1} & \text{continuous,}\\
          1 \cdot H_w^{-1}               & \text{binary.}
        \end{cases}
\end{equation}

% -----------------------------------------------------------------------
\section{Summary}

\begin{table}[h]
\centering
\renewcommand{\arraystretch}{1.4}
\begin{tabular}{lll}
\toprule
                     & \textbf{Continuous (OLS)} & \textbf{Binary (logistic)} \\
\midrule
Gradient $g_w$       & $\pm\bar{\xt} = \pm(1,\bar{x}^\top)^\top$
                     & $\pm n^{-1}\sum_i v_w(x_i)\xt_i$ \\
\texttt{VarCovMatrix} returned by \texttt{ds.glm}
                     & $(X_w^\top X_w)^{-1}$ (unscaled)
                     & $H_w^{-1}$ (full covariance) \\
Scale factor $\hsigma^2$
                     & $\sum_k(\mathrm{RSS}_k^{(1)}+\mathrm{RSS}_k^{(0)})/(n-2(d{+}1))$
                     & $1$ (fixed) \\
$V_w$ used in \eqref{eq:var}
                     & $\hsigma^2(X_w^\top X_w)^{-1}$
                     & $H_w^{-1}$ \\
Extra server disclosure
                     & $\mathrm{RSS}_k$ per site, per arm
                     & $\sum_i v_w(x_i)\xt_i$ per site, per arm \\
\bottomrule
\end{tabular}
\caption{Comparison of continuous and binary variance estimation paths.}
\end{table}

Both paths feed into the single formula~\eqref{eq:var}, which is implemented in
\texttt{.compute\_variance\_delta()} on the client side.

\end{document}
